% SPDX-FileCopyrightText: 2026 Yiyu Zhou <yiyu@yiyuzhou.io>
%
% SPDX-License-Identifier: GPL-3.0-or-later

\chapter{Running a show}
\label{ch:running}

cinq is a cue-based show controller in the spirit of QLab.  A show is a list
of cues, each of which does one thing, such as playing a sound or fading it
out.  The operator runs the show by pressing GO, which fires the next cue.

\section{Getting the show}

This tutorial comes with a storm show, in the \texttt{demo} directory beside
this PDF, such as \texttt{/usr/share/doc/cinq/demo}, along with the sounds,
picture, and film that the show plays.

\begin{enumerate}
\item Copy the \texttt{demo} directory to a place of your own, such as your
  home directory.  Chapter~\ref{ch:cues} saves a show of your own there,
  which the installed directory does not allow.
\item Start cinq, press \key{Control+o}, and open \texttt{storm.cinq.toml}
  in your copy.
\end{enumerate}

cinq's README explains how to install cinq.  Before the first GO, make sure
that the computer's sound is on.

\section{The first GO}

The main window opens as shown in Figure~\ref{fig:window}.  The transport at
the top names the cue on standby, which the next GO fires: cue~1, House open,
with its notes for the operator.  Below it is the cue list, and to its right
the sidebar shows the show's cue lists, its active cues, and any warnings.

\begin{figure}[H]
  \centering
  \includegraphics[width=\textwidth]{screenshots/window.png}
  \caption{Main window as the storm show opens}
  \label{fig:window}
\end{figure}

\begin{enumerate}
\item Press \ui{GO} or \key{Space}.  The preshow music starts, and the
  sidebar lists it under \ui{Active} with a progress bar and a \ui{Stop}
  button, as Figure~\ref{fig:playing} shows.
\item Press \ui{Panic} or \key{Escape}.  The music fades out over one second
  and stops.
\end{enumerate}

\begin{figure}[H]
  \centering
  \includegraphics[width=\textwidth]{screenshots/playing.png}
  \caption{Main window after the first GO, with the preshow music playing}
  \label{fig:playing}
\end{figure}

Cue~1 is a memo cue, which does nothing itself; it continues immediately to
cue~2, which plays the music, as Chapter~\ref{ch:sequences} explains.  After
each GO, standby moves past the cues that the GO fired, so cue~3 is now on
standby.

Panic fades out every playing cue over one second and stops every other
active cue.  A second panic during the fade stops everything at once.  The
keys for GO and Panic work anywhere in cinq, even while you type in a field.

\section{Standby}

GO fires the cue on standby, which is not necessarily the selected cue.

\begin{enumerate}
\item Click cue~7, Thunder crack, or select it with \key{ArrowDown}.  The
  inspector below the list shows its settings, but the transport still names
  cue~3.
\item Press \ui{Standby selected} or \key{Enter}.  The transport now names
  cue~7, and the \texttt{State} column marks it \texttt{NEXT}\@.
\item Press \ui{GO}\@.  The thunder crack plays.
\end{enumerate}

Selecting a cue changes only what the inspector shows, so the operator can
examine any cue without changing what GO fires.  \ui{Pause all} pauses every
running cue, and \ui{Resume all} resumes them.

\section*{Exercise}

You can now run a show: GO fires the cue on standby, Panic stops everything,
and \ui{Standby selected} chooses the next cue.  Fire cue~4, Distant thunder,
then put cue~1 back on standby to restart the show.
